% R. Brown's translations of two letters to L. Breen, with his note and comments (cote 134-2): archive cote 134-2, site page 137 (PDF page 138).
% translation of the letter of 17-19 July 1975. Draft compared: ps1.tex, lines 1572-1612.
% AI-transcribed by Claude Opus 5.5 (claude-opus-5-5, adaptive thinking, effort high), Message Batches API, batch brown_full (msgbatch_01Md3i5yK1PxkpJn5GHfEExE), 2026-11-13; not yet checked by eye.
\scanpage{134-2/137}
\mspage{- 10 -}

as an amalgamated sum.  In terms of this diagram of spaces,
the above functor $\phi$ interprets immediately as
\[ \phi(F_1) \simeq p_* j^*(\tilde{F}_1) \]
where $F_1 \to \tilde{F}_1$ is the restriction from $X_1$ to $\tilde{X}_1$ (which is
an equivalence of categories for the envisaged (locally constant)
sheaves).  Giving $F = (F_0 , F_1 , u : F_0 \to \phi F_1)$ can then also
be made explicit as giving
\[ F_0 , \tilde{F}_1 , \tilde{u} : p^*(F_0) \to j^*(\tilde{F}_1) \]
where $F_0(\tilde{F}_1)$ are locally constant sheaves on $X_0$ (respectively
$X_1$).  It is necessary to recall that up to here $p$ is a real
fibration, and $j$ is an inclusion (in practice, - for the case
$X_0$ , $X_1$ smooth - the inclusion of the boundary is a manifold
with boundary).

If you prefer, one can also take the diagram which is less
pretty (but a little more canonical)
\diagram{\begin{tikzcd}
\mathring{T} \ar[r,"j'"] \ar[d,swap,"p'"] & X_1 \\
X_0 &
\end{tikzcd}}
coming essentially to the same thing, as formed from spaces
homotopic to the preceding one.  One can even replace $X_0$ by $T$
($X_0$ being a deformation retract of it) and write
\diagram{\begin{tikzcd}
\mathring{T} \ar[r] \ar[d,swap,"p''"] & X_1 \\
T &
\end{tikzcd}}

\margin{138}
